% Parent-ready research fragment. Labels use eg:; no preamble is included.
% References to define in the final article: OrderedIncoming, GapSixIncoming,
% DLMFStirling. Sources inspected: current ProveIt incoming archives dated
% 2026-10-11 and 2026-10-12. The complete gap-six result is inherited.

\section{Effective reconstruction at every fixed gap}
\label{eg:section}

The incoming ordered-boundary report \cite{OrderedIncoming} constructs a Gamma integral basis
for the unequal-even-tail moment problem and proves eventual
nonsingularity for each fixed gap.  The subsequent gap-six report \cite{GapSix}
makes one of those eventual statements effective and completes that
ray by finite-field arithmetic.  We now give a fully explicit threshold
for \emph{every} fixed gap.  We then sharpen the estimates at gap eight
and certify all remaining indices.  In particular, the complete
gap-eight family, which the inspected gap-six report explicitly leaves
open, is resolved below.

The conclusions are statements about rational matrices on specified
formal labels.  They do not imply algebraic or rational independence
of the evaluated zeta values.  The contribution relative to the
inspected sources is the explicit general threshold and the completed
gap-eight ray; no historical-priority claim is made.

\subsection{The inherited rational obstruction}

Fix integers $m\ge1$ and $b>m$, and put
\[
 p=2b-2m,\qquad q=2b,\qquad x_n=n+\tfrac12,
 \qquad A_v(n)=\zeta(v,n+1).
\]
The meromorphic continuation of
\[
 D_{p,q}(s)=\sum_{n=0}^{\infty}A_p(n)A_q(n)x_n^{-s}
\]
is regular at every nonpositive even argument.  Indeed, the centered
Hurwitz expansion has only powers $x_n^{1-v-2h}$, so the possible
poles of the continued product sum are at the odd integers
$3-p-q-2h$.  Write $D_r=D_{p,q}(-2r)$.

The block coordinates from \cite{OrderedIncoming,GapSix} satisfy
\begin{align}
 D_r&=\sum_{j=0}^r\frac{\binom{2r}{2j}}{2j+1}
             B_{2r-2j}(1/2)C_j,\label{eg:forward}\\
 C_j&=\sum_{r=0}^j\binom{2j+1}{2r}2^{-2(j-r)}D_r.
 \label{eg:inverse}
\end{align}
For $j\ge b$ they have the late coordinates
\[
 C_j=\sum_{k=0}^{b-1}a_{jk}\zeta(2k),\qquad \zeta(0)=-\tfrac12.
\]
Put $d=2j+1$ and $e=d-p-q+2k+1$.  The coefficient is zero
if $e<0$; otherwise it is
\begin{equation}\label{eg:late-entry}
 a_{jk}=B_e\left\{
       \frac{\binom{d-q+1}{p-2k}}{d-q+1}
      +\frac{\binom{d-p+1}{q-2k}}{d-p+1}\right\}.
\end{equation}
Each summand is omitted unless its lower binomial index is between
$2$ and the upper index.  In particular the lower index zero is
omitted.  Define
\[
 M_{b,m}=(a_{b+r,k})_{0\le r,k<b}.
\]
The first $2b$ moment rows are expressed on the formal labels
\[
 C_0,\ldots,C_{b-1},\quad\zeta(0),\ldots,\zeta(2b-2).
\]
Their exact coefficient determinant is $\det M_{b,m}/(4b-1)!!$.
Thus nonsingularity of $M_{b,m}$ is exactly the optimal-prefix question.

We recall the fixed-gap reduction in a notation which keeps the
Bernoulli coefficients distinct from recurrence factors.  Set
\begin{align*}
 \beta_n&=\frac{2|B_{2n}|}{(2n)!},\qquad n\ge1,\\
 a_u&=0\quad(1\le u\le m),\qquad
 a_u=(2u)^{\underline{2m}}\quad(u>m),\\
 b_s&=(2s-1)^{\underline{2m}},\qquad
 d_s=2(a_s+b_s)\quad(s>m).
\end{align*}
Here the underline means falling factorial.  Define the rational
fundamental array by
\begin{align}
 R_{u,j}&=\delta_{uj}&&(1\le u,j\le m),\notag\\
 R_{s,j}&=\frac1{d_s}\sum_{u=1}^{s-1}
      \beta_{s-u}(b_s+a_u)R_{u,j}\quad&&(s>m).
 \label{eg:recurrence}
\end{align}
Finally let
\[
 F_{b,m}=(R_{b+i,j})_{1\le i,j\le m},\qquad f_{b,m}=\det F_{b,m}.
\]
The inherited complementary-minor identity is
\begin{equation}\label{eg:rational-bridge}
 \det M_{b,m}=
 \frac{(-1)^{b(b-1)/2}}{4^bb!}
       \left(\prod_{s=m+1}^{b+m}d_s\right) f_{b,m}.
\end{equation}
Every factor in the product is strictly positive.  Its normalization
includes the omitted-index convention in \eqref{eg:late-entry}.
This identity follows by coefficient comparison in the null equation
\[
 (R^{(2m)}-R^{(2m)}(0))\cot(z/2)
                         +(R\cot(z/2))^{(2m)}=0
\]
and the complementary-minor identity for its lower-triangular
coefficient operator; these inherited steps are reproduced in the
ordered-boundary and gap-six sources.

\subsection{Gamma realization and its signs}

Put
\begin{align*}
 W_m(x)&=\left|\Gamma\!\left(\tfrac14+\tfrac{ix}{2}\right)\right|^{4m},\\
 P_j(x)&=\sinh(\pi(m-j)x)\quad(1\le j<m),\qquad
 P_m(x)=\tanh(\pi x),\\
 \lambda_j&=j\pi,\qquad A_m=2^m(2\pi)^{2m},\qquad
 a_j^*=\begin{cases}A_m/2,&j<m,\\ A_m,&j=m.
                   \end{cases}
\end{align*}
The incoming Gamma-basis theorem \cite{OrderedIncoming} gives
\begin{equation}\label{eg:gamma-coefficients}
 t_{n,j}=\frac4{(2n)!}\operatorname{Im}
       \int_0^\infty(x+i/2)^{2n}P_j(x)W_m(x)\,dx
\end{equation}
and
\begin{equation}\label{eg:gamma-bridge}
 (t_{b+i,j})_{i,j=1}^m=F_{b,m}T_m,\qquad
 T_m=(t_{n,j})_{n,j=1}^m,\qquad
 \operatorname{sgn}\det T_m=(-1)^{m(m-1)/2}.
\end{equation}
All integrals converge absolutely.  The null equation follows from
the Gamma recurrence after two Fourier contours are shifted by
$\pm i$; the sinh modes cross no poles, and the tanh mode supplies
the two residues at $\pm i/2$ which give the subtracted derivative.
The sign in \eqref{eg:gamma-bridge} follows from continuous
Cauchy--Binet: its integrand is a positive weight times the product
of a Vandermonde in $x_r^2$ and a reversed Vandermonde in
$\cosh(\pi x_r)$.  These exact inputs, rather than numerical
normalizations of the Gamma integrals, will be used below.

\subsection{An explicit real-density estimate for arbitrary gaps}

\begin{lemma}\label{eg:density}
For $X=2m$, every $x\ge X$, and $1\le j\le m$,
\begin{equation}\label{eg:density-bound}
 \frac{P_j(x)W_m(x)}{a_j^*}
   =x^{-m}e^{-\lambda_jx}(1+\delta_{m,j}(x)),
 \qquad |\delta_{m,j}(x)|\le\frac{2m}{x^2}.
\end{equation}
On $0\le x\le X$ one has the elementary bound
\begin{equation}\label{eg:compact-density}
 0\le P_j(x)W_m(x)/a_j^*<2^{16m^2+12m}.
\end{equation}
\end{lemma}

\begin{proof}
The complex Stirling remainder \cite[\S5.11(ii)]{DLMF}, after retaining $1/(12z)$, gives
\[
 \log\Gamma(z)=(z-\tfrac12)\log z-z+\tfrac12\log(2\pi)
                        +\frac1{12z}+\rho(z),\qquad
 |\rho(z)|\le\frac1{90|z|^3}
\]
for $z=1/4+ix/2$.  Indeed the first omitted term has magnitude
$1/(360|z|^3)$ and the secant factor is at most four in this
closed right-half-plane sector.  Let $t=1/(2x)$ and
\[
 L_m(x)=\log\frac{W_m(x)}{A_mx^{-m}e^{-m\pi x}}.
\]
Taking $4m$ times the real part gives exactly
\[
 L_m(x)=-\frac m2\log(1+t^2)
      +m\left(\frac{\arctan t}{t}-1\right)
      +\frac{4m}{3}\frac{t^2}{1+t^2}+4m\operatorname{Re}\rho(z).
\]
Using $\log(1+t^2)\le t^2$ and
$0\le1-\arctan(t)/t\le t^2/3$, for $x\ge2$ we obtain
\begin{equation}\label{eg:log-density}
 |L_m(x)|\le\frac{13m}{24x^2}+\frac{16m}{45x^3}
          \le\frac{259m}{360x^2}<\frac{3m}{4x^2}.
\end{equation}
For $x\ge2m$ this is at most $3/(16m)<1$.  The elementary
inequality $|e^u-1|\le2|u|$ for real $|u|\le1$ now gives
an error at most $3m/(2x^2)$ for the normalized Gamma density.
The remaining normalized sinh or tanh multiplier lies in $[0,1]$
and differs from one by at most $2e^{-2\pi x}\le1/(9x^2)$.
Combining these estimates proves \eqref{eg:density-bound}.

Euler's integral gives $|\Gamma(1/4+ix/2)|\le\Gamma(1/4)<5$.
For $0\le x\le2m$, each multiplier is at most
$e^{2\pi m(m-1)}$.  Since $a_j^*>1$, $\pi<4$, $e<4$,
and $5<8$, the stated generous power-of-two bound follows.
\end{proof}

\subsection{An effective shifted Laplace lemma}

For $N>m+1$ and $\lambda>0$ define
\[
 S_N^{(r)}(\lambda)=
       \frac{\Gamma(N-m-r+1)}{\lambda^{N-m-r+1}},
       \qquad r=0,1,2.
\]
Let
\begin{equation}\label{eg:constants}
 C_m=\left\lceil4m+\frac{3m(m+1)}8+1\right\rceil,
 \qquad N_0(m)=4096(m+1)^2.
\end{equation}

\begin{lemma}\label{eg:laplace}
For every integer $N\ge N_0(m)$ and $1\le j\le m$,
\begin{equation}\label{eg:laplace-bound}
 \left|\frac1{a_j^*}\int_0^\infty(x+i/2)^NP_j(x)W_m(x)\,dx
   -e^{i\lambda_j/2}\left(S_N^{(0)}(\lambda_j)
                 +\frac{im}{2}S_N^{(1)}(\lambda_j)\right)\right|
       < C_m S_N^{(2)}(\lambda_j).
\end{equation}
\end{lemma}

\begin{proof}
Write $X=2m$.  On $[X,\infty)$ replace the normalized density
by $x^{-m}e^{-\lambda x}$.  Its absolute error is at most
\[
 2m\int_X^\infty |x+i/2|^Nx^{-m-2}e^{-\lambda x}\,dx.
\]
On $x\ge\sqrt N$ one has
$|x+i/2|^N\le x^N\exp(N/(8x^2))<2x^N$, so this part is
at most $4mS_N^{(2)}$.  On $X\le x\le\sqrt N$, the factor
$2mx^{-m-2}$ is at most one; hence the remaining part is at most
$\sqrt N(N+1)^{N/2}$.

For the model integral set $y=x+i/2$ and move the ray to
$[X,\infty)$.  Its integrand, after extracting $e^{i\lambda/2}$,
is $y^{N-m}(1-i/(2y))^{-m}e^{-\lambda y}$.  No pole is crossed.
On real $y\ge2m$, Taylor's theorem applied to $(1-u)^{-m}$ gives
\[
 \left|(1-i/(2y))^{-m}-1-\frac{im}{2y}\right|
       \le\frac{3m(m+1)}{8y^2}.
\]
For detail, $u=1/(2y)\le1/(4m)$, and the positive-coefficient
Taylor remainder is bounded by
$m(m+1)u^2(1-u)^{-m-2}/2$.  Also
\[
 (1-u)^{-m-2}\le
  \exp\!\left(\frac{(m+2)u}{1-u}\right)
  \le\exp\!\left(\frac{m+2}{4m-1}\right)<3.
\]
Thus the model remainder contributes at most
$3m(m+1)S_N^{(2)}/8$.

The short vertical segment at $X$, the two retained integrals
between $0$ and $X$, and the original compact density integral
are together bounded by $2^{20(m+1)^2}(3m)^N$.  Indeed the
compact density bound contributes at most
$X2^{16m^2+12m}(X+1)^N$; the other three terms are bounded
by $(2+m/2)(X+1)^N$.  The displayed larger constant covers their
sum for every $m\ge1$.

It remains to dominate
\[
 R_N=2^{20(m+1)^2}(3m)^N+\sqrt N(N+1)^{N/2}
\]
by $S_N^{(2)}$.  Put $k=N-m-2$.  Since $\lambda_j<4m$ and
$k!\ge(k/3)^k$,
\[
 S_N^{(2)}(\lambda_j)>
             \frac1{4m}\left(\frac{k}{12m}\right)^k.
\]
For $N\ge N_0(m)$, one has $k\ge3N/4$ and
$\sqrt{N+1}\le(5/4)\sqrt N$.  Consequently
\begin{align*}
 \frac{\sqrt N(N+1)^{N/2}}{S_N^{(2)}}
  &\le N^{m+4}\left(\frac{20m}{\sqrt N}\right)^N
    <N^{m+4}2^{-N},\\
 \frac{2^{20(m+1)^2}(3m)^N}{S_N^{(2)}}
  &\le2^{20(m+1)^2}N^{m+3}
                \left(\frac{48m^2}{N}\right)^N
    <2^{20(m+1)^2}N^{m+3}4^{-N}.
\end{align*}
The elementary inequality $N^{m+4}\le2^{N-2}$ holds throughout
this range.  At $N_0(m)$ it follows from
$\log_2N_0(m)\le12+2m$ and
$(m+4)(12+2m)<N_0(m)-2$; it persists by induction because
$N>2(m+4)$ implies $(1+1/N)^{m+4}<2$.
The two ratios are therefore less than $1/4$ each, the second
using also $N>20(m+1)^2$.  Thus $R_N<S_N^{(2)}/2$.
Adding all errors proves \eqref{eg:laplace-bound}.
\end{proof}

\subsection{A closed, computable threshold for every fixed gap}

Let $V_m=(j^{-2i})_{1\le i,j\le m}$ and define the positive rational
number
\begin{equation}\label{eg:delta}
 \delta_m=\frac{|\det V_m|}{\operatorname{perm}V_m}.
\end{equation}
Here $\operatorname{perm}V_m$ is the sum of its $m!$ unsigned
permutation products.  Thus $0<\delta_m\le1$, and its numerator
is explicitly the absolute reciprocal-square Vandermonde product
\[
 |\det V_m|=
 \left(\prod_{j=1}^m j^{-2}\right)
             \prod_{j<k}(j^{-2}-k^{-2}).
\]
Set
\begin{equation}\label{eg:B}
 E_m=8C_m+2(m-1),\qquad
 \boxed{\displaystyle
 B(m)=\max\left\{2048(m+1)^2,
       1+\left\lceil\frac{2mE_m}{\delta_m}\right\rceil\right\}.}
\end{equation}

\begin{theorem}[Effective reconstruction on every fixed-gap ray]
\label{eg:effective}
For every integer $m\ge1$ and every $b\ge B(m)$,
\begin{equation}\label{eg:general-sign}
 f_{b,m}\ne0,\qquad
 \operatorname{sgn}f_{b,m}=(-1)^{\lfloor m^2/4\rfloor}.
\end{equation}
Consequently the first $q=2b$ moments for
$(p,q)=(2b-2m,2b)$ form an optimal formal prefix, and
\[
 \operatorname{sgn}\det M_{b,m}
       =(-1)^{b(b-1)/2+\lfloor m^2/4\rfloor}.
\]
For any prescribed $m$, all possible exceptions lie in the explicitly
specified finite interval $m+1\le b<B(m)$.
\end{theorem}

\begin{proof}
Let $N_i=2b+2i$ for $1\le i\le m$.  At odd $j$, the leading
phase in \eqref{eg:laplace-bound} is $\pm i$, and at even $j$
it is $\pm1$.  Thus the odd-mode leading term is
\[
 \sigma_j\frac{4a_j^*}{N_i!}S_{N_i}^{(0)}(\lambda_j),
 \qquad \sigma_j=\sin(j\pi/2),
\]
with relative error bounded by
\[
 \frac{C_m\lambda_j^2}{(N_i-m)(N_i-m-1)}
       <\frac{16C_mm^2}{(N_i-m-1)^2}.
\]
At even $j$ the first nonzero imaginary term is
\[
 \sigma_j\frac{2ma_j^*}{N_i!}S_{N_i}^{(1)}(\lambda_j),
 \qquad \sigma_j=\cos(j\pi/2),
\]
with relative error less than $8C_m/(N_i-m-1)$.
Retaining this second phase term is essential.

Extract the positive row factors
\[
 \frac{4\Gamma(N_i-m+1)}{\Gamma(N_i+1)}\pi^{-2i}
\]
and the positive column factors $a_j^*\lambda_j^{-2b+m-1}$,
multiplying each even column factor additionally by
$m\lambda_j/(2(2b+2-m))$.  Also extract its sign $\sigma_j$.
The remaining array $E=(E_{ij})$ has the form
\[
 E_{ij}=j^{-2i}(1+\eta_{ij}).
\]
For an even column the extra row factor is
\[
 \rho_i=\frac{2b+2-m}{2b+2i-m},\qquad
 1-\rho_i\le\frac{2(m-1)}{2b+2-m}.
\]
Since $b\ge B(m)>m$ and $b\ge2m^2$, the preceding estimates give
\begin{equation}\label{eg:eta}
 |\eta_{ij}|<\frac{E_m}{b}<\frac{\delta_m}{2m}.
\end{equation}
For the odd columns use
$16C_mm^2/b^2\le8C_m/b$; for the even columns use
$|\rho_i(1+\epsilon)-1|\le1-\rho_i+|\epsilon|$.

Expanding the determinant into its signed permutation products,
\[
 |\det E-\det V_m|
 <\operatorname{perm}V_m
       \left\{\left(1+\frac{\delta_m}{2m}\right)^m-1\right\}
 <\operatorname{perm}V_m\,\delta_m=|\det V_m|.
\]
The last strict inequality holds for $0<\delta_m\le1$: the binomial
sum is bounded above by the corresponding infinite geometric sum
$\delta_m/(2-\delta_m)\le\delta_m$, with a strictly positive
omitted tail.  Hence $\det E$ has the sign of the reversed
Vandermonde, namely $(-1)^{m(m-1)/2}$.
The extracted column signs have product
$\prod_j\sigma_j=(-1)^{\lfloor m^2/4\rfloor}$.
Dividing by the initial determinant in \eqref{eg:gamma-bridge}
cancels the reversed-Vandermonde sign and proves
\eqref{eg:general-sign}.  Apply \eqref{eg:rational-bridge} to
obtain the assertion about the original moment matrix.
\end{proof}

The formula is deliberately conservative and does not claim a least
threshold.  Its first values are
\[
 \begin{array}{c|rrrrrr}
 m&1&2&3&4&5&6\\ \hline
 B(m)&8192&18432&32768&51200&635656&16146233.
 \end{array}
\]
The earlier complete gap-two, gap-four, and gap-six theorems are
stronger on their individual rays.  The advance here is an effective
formula valid for arbitrary prescribed $m$.  It does not assert
that every one of the finite remaining matrices is nonsingular;
the unrestricted all-pairs conjecture remains open.

\section{The complete gap-eight theorem}
\label{eg:gap-eight}

For $m=4$ sharper constants reduce the finite-completion threshold
from $51200$ to $8192$.  The next two estimates supply that reduction.

\begin{lemma}[Gap-eight density and phase estimates]
\label{eg:eight-laplace}
For $m=4$, $x\ge2$, and $1\le j\le4$,
\[
 \frac{P_j(x)W_4(x)}{a_j^*}
    =x^{-4}e^{-\lambda_jx}(1+\delta_j(x)),\qquad
 |\delta_j(x)|<\frac5{x^2}.
\]
For every integer $N\ge16384$,
\begin{equation}\label{eg:eight-phase}
 \left|\frac1{a_j^*}\int_0^\infty(x+i/2)^NP_j(x)W_4(x)\,dx
 -e^{i\lambda_j/2}
     \left(\frac{\Gamma(N-3)}{\lambda_j^{N-3}}
            +2i\frac{\Gamma(N-4)}{\lambda_j^{N-4}}\right)\right|
       <16\frac{\Gamma(N-5)}{\lambda_j^{N-5}}.
\end{equation}
\end{lemma}

\begin{proof}
Equation \eqref{eg:log-density} gives $|L_4(x)|<3/x^2\le3/4$.
The elementary Taylor-series bound $e^{3/4}<17/8$ implies
$|e^u-1|\le(3/2)|u|$ for real $|u|\le3/4$.  Adding the
normalized multiplier error $1/(9x^2)$ gives
\[
 |\delta_j(x)|<\left(\frac92+\frac19\right)x^{-2}<5x^{-2}.
\]
On $0\le x\le2$ the normalized density is less than $2^{80}$,
since $5^{16}<2^{38}$, $e^{6\pi}<3^{24}<2^{39}$, and $a_j^*>1$.

Repeat the proof of Lemma~\ref{eg:laplace} with $X=2$.  The
density error above $\sqrt N$ contributes less than
$10\Gamma(N-5)/\lambda_j^{N-5}$.  The model Taylor remainder
has the sharper bound
\[
 \left|(1-i/(2y))^{-4}-1-\frac{2i}{y}\right|<\frac5{y^2}
       \qquad(y\ge2).
\]
Indeed the positive series remainder divided by $u^2$, where
$u=1/(2y)\le1/4$, is increasing and equals $1504/81$ at $u=1/4$;
its contribution is therefore $376/(81y^2)<5/y^2$.

All compact and lower-range contributions are at most
\[
 R_N=2^{82}3^N+\sqrt N(N+1)^{N/2}.
\]
Now $\lambda_j\le4\pi<13$ and, with $k=N-6$,
\[
 \frac{\Gamma(N-5)}{\lambda_j^{N-5}}
       >\frac1{13}\left(\frac{k}{39}\right)^k.
\]
For every $N\ge16384$,
\[
 \frac{39\sqrt{N+1}}{N-6}<\frac12,\qquad
 \frac{117}{N-6}<\frac1{128},\qquad N^7\le2^{N-2}.
\]
The first inequality follows by squaring, checking the resulting
quadratic at $16384$, and noting its positive derivative thereafter;
the second is immediate; the third holds at $16384$ and persists
by the same induction as in Lemma~\ref{eg:laplace}.  Therefore
the two normalized parts of $R_N$ are bounded, respectively, by
$2^{82}N^6 2^{-7N}<1/4$ and $N^7 2^{-N}\le1/4$.
The total error is consequently less than $(10+5+1/2)$ times the
displayed Gamma factor, proving \eqref{eg:eight-phase}.
\end{proof}

\begin{proposition}[Certified positive tail at gap eight]
\label{eg:eight-positive}
For every integer $b\ge8192$, $f_{b,4}>0$.
\end{proposition}

\begin{proof}
In \eqref{eg:eight-phase}, an odd column has relative error less
than $2704/(N-5)^2$, and an even column has relative error less
than $104/(N-5)$; these use only $\lambda_j<13$.
Perform the row and column normalizations in Theorem~\ref{eg:effective}.
The even-column row factors are now
$\rho_i=(2b-2)/(2b+2i-4)$, with
$1-\rho_i\le6/(2b-2)$.
Hence the normalized entries satisfy
\[
 E_{ij}=j^{-2i}(1+\eta_{ij}),\qquad
 |\eta_{ij}|<\frac1{128}\qquad(b\ge8192).
\]
For the even columns use
$6/(2b-2)+104/(2b-3)<110/(2b-3)<1/128$;
the odd-column inequality is weaker still.

The four-by-four limiting matrix has exact values
\[
 \det V_4=\frac{175}{127401984},\qquad
 \operatorname{perm}V_4=\frac{157183}{4586471424},\qquad
 \frac{\det V_4}{\operatorname{perm}V_4}=\frac{6300}{157183}.
\]
Moreover
\[
 \det V_4-\operatorname{perm}V_4
        \left\{(129/128)^4-1\right\}
       =\frac{357064449025}{1231171548132409344}>0.
\]
Thus $\det E>0$.  The extracted signs are $(+,-,-,+)$, whose
product is positive, and $\det T_4>0$.  Equation
\eqref{eg:gamma-bridge} proves $f_{b,4}>0$.
\end{proof}

\subsection{An exact finite-field completion}

The cotangent coefficients admit the elementary recurrence
\begin{equation}\label{eg:beta-recurrence}
 \beta_1=\frac16,\qquad
 \beta_n=\frac1{2(2n+1)}\sum_{k=1}^{n-1}\beta_k\beta_{n-k}
                  \quad(n\ge2).
\end{equation}
It follows by substituting the Taylor expansion of
$(z/2)\cot(z/2)$ into its Riccati equation.
Equations \eqref{eg:beta-recurrence} and \eqref{eg:recurrence}
can be computed exactly modulo the prime
$\ell=2147483647=2^{31}-1$.

For $m=4$ and $5\le b\le8191$ the largest fundamental row
required is $8195$, so only $\beta_1,\ldots,\beta_{8194}$ are needed.
Every denominator is invertible modulo $\ell$.  The denominators
$2(2n+1)$ have their prime factors below $\ell$, and
\begin{equation}\label{eg:diagonal-factorization}
 d_s=2\frac{(2s)!}{(2s-8)!}\frac{2s-4}{s}\qquad(s\ge5)
\end{equation}
has the same property throughout the stated finite range.  Thus
induction identifies every computed residue with the reduction of
the exact rational entry.

\begin{proposition}[Finite gap-eight certificate]
\label{eg:eight-finite}
For every integer $5\le b\le8191$, the rational determinant
$f_{b,4}$ has a nonzero reduction modulo $2147483647$.
\end{proposition}

\begin{proof}[Exact computational proof]
The accompanying \texttt{verify\_gap\_eight.cpp} uses integer
arithmetic only.  It verifies primality by trial division through
$\lfloor\sqrt\ell\rfloor$, generates \eqref{eg:beta-recurrence},
generates the four fundamental columns from \eqref{eg:recurrence},
and evaluates each four-by-four determinant by finite-field Gaussian
elimination.  All $8187$ determinant residues are nonzero; the full
indexed list is supplied in \texttt{gap\_eight\_certificate.json}.
As an independent exact check of the cotangent normalization, the
program verifies all $8194$ coefficients by multiplying
$1-\tfrac12\sum_{n\ge1}\beta_nt^n$ by the directly generated
series of $\sin(\sqrt t/2)/(\sqrt t/2)$ and comparing with the
direct series of $\cos(\sqrt t/2)$.
The separate SymPy verifier builds rational Bernoulli numbers from
their standard definition and checks the original matrix
\eqref{eg:late-entry}, the full bridge \eqref{eg:rational-bridge},
and the finite-field residues independently for $5\le b\le16$.
These latter checks guard against a sign, factorial, or omitted-index
error in passing from the source matrix to the recurrence.
\end{proof}

\begin{theorem}[Every gap-eight pair has the optimal prefix]
\label{eg:eight-main}
For every integer $b\ge5$, the matrix $M_{b,4}$ is nonsingular.
For $p=2b-8$ and $q=2b$, consequently,
\begin{align*}
 \operatorname{span}_{\mathbb Q}\{D_0,D_1,\ldots\}
 &=\operatorname{span}_{\mathbb Q}\{D_0,\ldots,D_{q-1}\}\\
 &=\operatorname{span}_{\mathbb Q}
       \{C_0,\ldots,C_{b-1},\zeta(0),\ldots,\zeta(2b-2)\}.
\end{align*}
The first $q$ moment rows have formal rank $q$ and coefficient
determinant $\det M_{b,4}/(4b-1)!!$.  This is optimal for those
$q$ formal coordinates.
\end{theorem}

\begin{proof}
Proposition~\ref{eg:eight-positive} covers $b\ge8192$, and
Proposition~\ref{eg:eight-finite} covers $5\le b\le8191$.
Equation \eqref{eg:rational-bridge} proves nonsingularity of the
source matrix.  The triangular transforms
\eqref{eg:forward}--\eqref{eg:inverse} give the stated generation
and formal-rank conclusions.
\end{proof}

The certificate proves nonvanishing at the finite indices; it does
not determine their real signs.  Direct rational checks give
$f_{b,4}>0$ for $5\le b\le16$, and the analytic theorem gives
positivity for $b\ge8192$.  We make no positivity assertion for
the intervening finite interval on the basis of the residue test.

\subsection{An all-order family of absolutely convergent identities}

The completed ray produces ordinary convergent evaluations of
polygamma-tail products at every even pair separated by eight.
Let $(\lambda_b,\ldots,\lambda_{2b-1})$ be row zero of
$M_{b,4}^{-1}$ and define
\begin{align*}
 K_b(y)&=\sum_{j=b}^{2b-1}\lambda_jy^{2j+1},\\
 P_b(x)&=K_b(x+\tfrac12)-K_b(x-\tfrac12)
             =\sum_{r=0}^{q-1}w_rx^{2r}.
\end{align*}
For $0\le h\le4$ put
\[
 c_{v,0}=\frac1{v-1},\qquad
 c_{v,h}=\frac{B_{2h}(1/2)}{(2h)!}(v)_{2h-1}\quad(h\ge1),
 \qquad u_h=\sum_{e=0}^h c_{p,e}c_{q,h-e},
\]
and set
\[
 S_b(x)=\sum_{h=0}^4
       \left(\sum_{e=0}^h u_e w_{q-1-h+e}\right)x^{8-2h}.
\]
Every coefficient here is rational and specified by finite algebra.

\begin{corollary}[Rational polygamma-product evaluations]
\label{eg:eight-identities}
For every $b\ge5$, $p=2b-8$, $q=2b$, and $x_n=n+1/2$,
\begin{equation}\label{eg:eight-sum}
 \boxed{\displaystyle
 \sum_{n=0}^{\infty}\left\{
 \frac{P_b(x_n)}{(p-1)!(q-1)!}
       \psi^{(p-1)}(n+1)\psi^{(q-1)}(n+1)-S_b(x_n)
                         \right\}=-\frac12.}
\end{equation}
The sum converges absolutely.  The polynomial $P_b$ has exact degree
$2q-2$ and leading coefficient
\begin{equation}\label{eg:leading}
 w_{q-1}=\frac{4b(4b-1)f_{b-1,4}}{d_{b+4}f_{b,4}}\ne0,
\end{equation}
where at the initial index one uses
\[
 f_{4,4}=\frac1{13105577681191934125041254400000000}.
\]
\end{corollary}

\begin{proof}
By construction, $\sum_{j=b}^{2b-1}\lambda_j C_j=\zeta(0)=-1/2$.
The centered difference defining $P_b$ has coefficients
\[
 w_r=\sum_{j=\max\{b,r\}}^{2b-1}
       \lambda_j\binom{2j+1}{2r}2^{-2(j-r)}.
\]
Thus \eqref{eg:inverse} gives $\sum_rw_rD_r=-1/2$.
The first five terms of the centered tail product are
\[
 A_p(n)A_q(n)=\sum_{h=0}^4u_hx_n^{2-p-q-2h}
                      +O(x_n^{-p-q-8}).
\]
Multiplication by $P_b$ leaves precisely $S_b$ as its nonnegative
powers, with the remainder $O(x_n^{-2})$.  The analytically
continued sums of $1,x_n^2,\ldots,x_n^8$ vanish because
$\zeta(-2h,1/2)=0$.  Therefore the ordinary absolutely convergent
sum in \eqref{eg:eight-sum} equals that moment combination.
Finally, for even $v$,
$A_v(n)=\psi^{(v-1)}(n+1)/(v-1)!$.

For $b\ge6$, deleting the last row and first column of $M_{b,4}$
and reindexing $j,k,p,q$ downward by $1,1,2,2$, respectively,
gives exactly $M_{b-1,4}$.  Cramer's rule and
\eqref{eg:rational-bridge} give
$\lambda_{2b-1}=4b f_{b-1,4}/(d_{b+4}f_{b,4})$.
The centered difference multiplies the top coefficient by $4b-1$,
proving \eqref{eg:leading}.  At $b=5$ the same relation is
verified directly with the displayed rational auxiliary determinant.
Theorem~\ref{eg:eight-main} proves all other required nonvanishing.
\end{proof}


\subsection{A concrete pair beyond the previous small-order coverage}

For $(p,q)=(14,22)$, namely $b=11$, the construction gives an even
polynomial $P_{11}$ of degree $42$ and a subtraction polynomial
$S_{11}$ of degree $8$. The leading coefficient is
\[
 [x^{42}]P_{11}(x)=
 \frac{452253373522332335630380}{10060146014319869101877031}.
\]
All $22$ coefficients of $P_{11}$ and all five coefficients of
$S_{11}$ are supplied as exact rationals in
\texttt{gap\_eight\_14\_22\_example.json}.
The short generator \texttt{create\_gap\_eight\_example.py} checks
the inverse row against the original Bernoulli matrix and verifies
the entire subtraction polynomial symbolically. Thus the example
specifies a directly usable ordinary sum, without printing several
pages of large rational coefficients. The finite definitions above
are the full exact formula for every other pair as well.
